% Generated draft: requires mathematical and source review.
\documentclass[11pt]{article}
\usepackage[margin=25mm]{geometry}
\usepackage{fontspec}
\setmainfont{Latin Modern Roman}
\usepackage{amsmath,amssymb,mathtools,unicode-math}
\setmathfont{Latin Modern Math}
\usepackage{enumitem,needspace,xurl,hyperref}
\setlength{\parindent}{0pt}
\setlength{\parskip}{5pt}
\setlength{\emergencystretch}{3em}
\newcommand{\N}{\mathbb{N}}
\newcommand{\Z}{\mathbb{Z}}
\newcommand{\Q}{\mathbb{Q}}
\newcommand{\R}{\mathbb{R}}
\newcommand{\C}{\mathbb{C}}
\newcommand{\F}{\mathbb{F}}
\newcommand{\A}{\mathbb{A}}
\newcommand{\PP}{\mathbb{P}}
\newcommand{\E}{\mathbb{E}}
\newcommand{\eps}{\varepsilon}
\newcommand{\dd}{\,\mathrm{d}}
\begin{document}
\section*{UnsolvedMath ID: 30005642}
\textbf{Source problem code:} catalogue entry
\subsection{Definitions and mathematical statement}
The supplied record refers to geometric \(\ell\)-adic local systems \(V_\ell\) on a space \(X\), but its excerpt begins with ``With the above assumptions'' and omits those assumptions. In particular, it does not specify the base field, the hypotheses on \(X\), the geometric construction of \(V_\ell\), the allowed primes, or the meaning of the point set denoted by \(|X|\). The following is therefore a conditional transcription of the intended questions; the omitted setup must be recovered before this becomes a standalone mathematical statement.

Retain the source notation \(|X|^{\mathrm{tor}}_{V_\ell}\) for its toric locus and set
\[
T_\ell:=|X|^{\mathrm{tor}}_{V_\ell}\cap |X|.
\]
The supplied excerpt does not define the pointwise group \(G_{\ell,x}\) or specify whether the toric condition concerns that whole group or its identity component. Accordingly, \(T_\ell\) here names the source's intended locus without supplying an unverified definition.

Subject to all the omitted assumptions of the source's Conjecture 2, its non-density assertion is
\[
\overline{T_\ell}^{\,\mathrm{Zar}}\ne X
\qquad\text{for every prime \(\ell\) allowed by that setup.}
\]
The associated independence question asks whether, for the systems arising from the same source setup,
\[
T_\ell=T_{\ell'}
\qquad\text{for every pair of allowed primes \(\ell,\ell'\).}
\]
The common point domain and the relation between the systems must be specified in the missing setup. The intersection with \(|X|\) is retained exactly as in the supplied clean statement; its role cannot be resolved from this excerpt.

The excerpt reports the following restricted result, which must be distinguished from the unrestricted conjecture. For a prime \(\ell\), its level condition is
\[
(\mathrm{Lev}_\ell):\quad
\operatorname{im}\!\left(\pi_1(X_{\bar k})\longrightarrow
\operatorname{GL}(V_\ell)\right)
\text{ is a pro-\(\ell\) group}.
\]
Here \(k\), \(\bar k\), the representation, and its fiber are part of the omitted setup. Assuming the hypotheses of Conjecture 2 and \((\mathrm{Lev}_{\ell_1})\), \((\mathrm{Lev}_{\ell_2})\) for two distinct primes \(\ell_1\ne\ell_2\), the extracted Theorem 1 asserts the existence of a set \(L\) of primes of positive Dirichlet density such that, for every \(\ell\in L\),
\[
\overline{T_\ell^{-}\cap X(k)}^{\,\mathrm{Zar}}\ne X.
\]
If the complex period map describing \(V_\infty\) is finite-to-one, it further asserts that \(T_\ell^{-}\cap X(k)\) is finite for these primes. The apparent reading of the damaged definition of the minus locus is
\[
T_\ell^{-}=\{x\in T_\ell:\text{ there is an allowed prime }\ell_x\ne\ell
\text{ with }x\in T_{\ell_x}\}.
\]
This reading requires source verification: superscripts and subscripts in the supplied extraction are corrupted. In particular, the reported theorem concerns a minus locus intersected with \(X(k)\), whereas the conjecture is stated using \(|X|\); these domains must not be identified without the missing assumptions.

The excerpt also reports that the solvable locus
\[
|X|^{\mathrm{solv}}_{V_\ell}
:=\{x\in|X|:G_{\ell,x}\text{ is solvable}\}
\]
is independent of \(\ell\) and is contained in an exceptional locus denoted \(|X|_{V_\ell}\), which the excerpt does not define. It attributes independence to class field theory, citing \([\mathrm{Se68}]\), and states that this holds more generally for any \(\Q_\ell\)-local system that is almost pointwise geometric in the sense of Fontaine--Mazur; that condition is not defined in the record.

The motivated toric locus, denoted \(|X|^{\mathrm{tor}}_{V_{\mathrm{mot}}}\), is described as the set of points \(x\in|X|\) for which the identity component of the motivated Galois group, in the sense of Andr\'e, of the motivated motive \(h^i(Y_x)(j)\) is a torus. The family \(Y\), the indices \(i,j\), and the required motivic setup are omitted. The excerpt asserts its inclusion in the toric locus for every prime \(\ell\), and reports, under the restricted theorem's hypotheses,
\[
\overline{|X|^{\mathrm{tor}}_{V_{\mathrm{mot}}}\cap X(k)}^{\,\mathrm{Zar}}\ne X.
\]
These are supplied background assertions, not independently verified results of this draft.

Finally, the excerpt says that Conjecture 2 follows from an unspecified Conjecture 1, and also from the Mumford--Tate and generalized Andr\'e--Oort conjectures through a deep implication involving the average Colmez conjecture. It associates the intended toric locus with CM points of a polarizable rational variation of Hodge structure \(V_\infty\), whose displayed construction and comparison are corrupted. It explicitly says that the general comparison with the CM locus, and even independence of \(\ell\), are unknown in that discussion. Its final remark expresses hope for an extension to arbitrary \(d\ge1\) without level conditions, still restricting \(\ell\) to a positive-density set of primes, but the remark is truncated and \(d\) is not defined.
\subsection{Short English statement}
Under the report's omitted geometric hypotheses, the conjecture asks whether the toric Galois locus is not Zariski dense in \(X\). A related question asks whether that locus is independent of \(\ell\). The supplied excerpt reports a restricted theorem involving level conditions at two distinct primes, a positive-density set of primes, and a smaller locus of \(k\)-rational points, with finiteness additionally requiring a finite-to-one period map. Missing hypotheses and damaged notation prevent a complete standalone formulation. This is a statement-only draft requiring human mathematical review.
\subsection{Sources}
\begin{itemize}
\item Local UnsolvedMath registry: supplied complete record for permanent ID \(30005642\), title \emph{Non-Density of Toric Galois Loci}, exact source problem code \texttt{catalogue entry}. The record supplies no individual proposer. Its proposed year is \(2023\).
\item Supplied report citation: \emph{MFO--RIMS Tandem Workshop: Arithmetic Homotopy and Galois Theory (2024)}, Oberwolfach Reports, pp.\ 2377--2488, DOI \texttt{10.4171/owr/2023/42}. Supplied source URL: \url{https://doi.org/10.4171/owr/2023/42}. The supplied title year \(2024\) and the record's \(2023\) year information are retained without reconciliation.
\item Supplied literature source, labelled ``Cadoret--Stix report'': \url{https://doi.org/10.4171/owr/2023/42}. The background describes Oberwolfach Reports as records of mathematical workshops and their open-problem discussions.
\item Literature assessment dated exactly \texttt{2026-08-22}: the supplied curation labels the problem ``partially solved'' and says that a theorem under level conditions at two primes gives non-density and conditional finiteness conclusions while the unrestricted conjecture remains open. This is supplied dated curation derived from the OWR record, not an independently verified assessment; its time-sensitive status and cited claims require independent checking.
\item The registry labels the statement \texttt{corrected\_verified} and reports that it was checked against the official report. That verification claim is supplied metadata and was not independently repeated for this draft.
\item The extracted discussion cites \([\mathrm{Se68}]\) for a class-field-theoretic assertion. No corresponding bibliographic entry is supplied, so this citation cannot be identified reliably from the record alone.
\end{itemize}
\end{document}
